% The twoside setting lets LaTeX know you want 2 different header or footer
% layouts. Not Essential but it looks really nice.
\documentclass[twoside, a4paper, 12pt]{article}

% Package Settings
% The UTF-8 Setting just allows full unicode chars to be used, especially
% handy for citing things from other languages.
\usepackage[utf8]{inputenc}
% Geometry just sets up the page to use proper margins. Check your subject's
% requirements for this, some want 2.5cm all around the doc.
\usepackage[margin=2cm, left=3cm, includefoot, asymmetric]{geometry}
% Using double spacing. Pretty simplistic really
\usepackage{setspace}

% 2017 Sem 1 History subjects require Chicago 16A Footnote Style referencing.
% I use the authordate-trad environment because it prints all of the
% authors in the footcite rather than just the first.
\usepackage[backend=biber,authordate-trad]{biblatex-chicago}
% All your references should go into Refs.bib. A sample reference layout
% Will be included, but you can google `bibtex citation generator' and
% Paste in a webpage or book name into the first result and get a decent
% bibtex style citation that can be used here.
\addbibresource{Refs.bib}
% This lets you include pictures! Do some research because it can be fiddly.
\usepackage{graphicx}
% This one makes for easy quotations and block quotations. Once again, do
% Research if you need.
\usepackage{csquotes}
% This one just lets you make fancy headers and footers with complex layouts.
\usepackage{fancyhdr}
% This is just here to let me put the title in the footer but it comes in handy
\usepackage{titling}
% This makes all of your URLs into hyperlinks and links your citations around
% the document.
\usepackage[hidelinks]{hyperref}
% Another environment for putting figures in, lets the words wrap around nicely
\usepackage{wrapfig}
% Allows for more useable captions for figures.
\usepackage{caption}
% Literally lets me use the lipsum command for junk text. Can delete
\usepackage{lipsum}

% Page Style Stuff
\fancypagestyle{plain}{
  \fancyhf{}
  %Change these values to change the thicknesses of the top and bottom lines!
  \renewcommand{\headrulewidth}{0.1pt}
  \renewcommand{\footrulewidth}{0.4pt}
  %Personalised footer
  \fancyfoot[LE,RO]{\theauthor}
  \fancyfoot[C]{\thetitle}
  \fancyfoot[RE,LO]{Page \thepage}
}

% Allows for sourcing images properly.
\newcommand{\source}[1]{\caption*{Source: {#1}} }

\title{The Title Goes Here}
\author{PhalanxHead}
\date{\today}

% Handy Citation Notes for later:
% \footcite[200]{morris_powell_2010} %To add page numbers to citations
% \blockquote{} % If you plan to quote a huge section of a book.

\begin{document}

\pagenumbering{Alph}
\begin{titlepage}
    \centering
    \includegraphics[width=0.3\textwidth]{UniLogo.jpg}\par\vspace{1cm}
    {\scshape\Large THE UNIVERSITY OF MELBOURNE \\
              \large Subject Code and Name \par} % The \\ splits the line
    \vspace{1.5cm}
    {\Huge \thetitle \par}
    \vfill

% Bottom of the page
    {\Large\itshape \theauthor \hspace{1em} -- \hspace{1em} STUDENT NUMBER \par}
    \vspace{1.5cm}
    {\Large \thedate}
\end{titlepage}
% Format first page correctly.
\pagestyle{plain}
\pagenumbering{arabic}

% Essay Section should go in the doublespacing environment
\begin{doublespacing}
This is the basic layout of your typical essay. {\LaTeX} looks really nice
and is definitely worth using for your essays. The citations alone are really
easy to use and take the stress out of actually writing the essay. \footcite{PhalanxHead_2017}
\par\vspace{1em}

\lipsum[1-5]

\end{doublespacing}

\nocite{*}
% Bib page
\newpage

\printbibliography

\end{document}
